\capituloSinNumero{Situación problemática}
\label{cap:situacion}

La situación problemática se concentra en la ineficiencia del proceso de
venta de paquetes turísticos en Vizcacha Tours, Potosí, que afecta
principalmente el control de los cupos disponibles y la confirmación de las
reservas. El árbol de problemas que sintetiza las causas y los efectos de
esta situación se presenta en la figura~\ref{fig:arbol}.

Entre las posibles causas están la falta de un criterio uniforme para revisar
la disponibilidad antes de confirmar, la ausencia de reglas claras para
modificar los estados de reserva y la demora en comunicar cambios de
capacidad o fechas. Estos factores no dependen necesariamente del uso de
registros manuales o digitales: pueden presentarse si el proceso no está bien
definido.

De confirmarse estas dificultades, podrían ocasionar reservas mayores a la
capacidad de una salida, reprogramaciones, tiempo adicional de revisión y
reclamos por parte de los turistas. Además, la información de cotizaciones,
pagos o proveedores podría quedar descoordinada respecto de las reservas
confirmadas.

\begin{landscape}
\thispagestyle{paisaje}
\begin{figuraapa}{fig:arbol}{Árbol de problemas del proceso de venta de paquetes turísticos}
\centering
\scalebox{0.85}{%% =====================================================================
%%  arbol-problemas.tex
%%  Diagrama de árbol de problemas (TikZ).
%%  Se incluye desde el Anexo A en una página horizontal.
%% =====================================================================
\begin{tikzpicture}[
  nodo/.style={
    draw=black!65, line width=0.5pt, rounded corners=1pt,
    align=center, font=\scriptsize, text width=4cm,
    inner sep=3pt, minimum height=1.7cm, fill=white
  },
  central/.style={
    nodo, text width=17cm, font=\small\bfseries,
    fill=black!8, minimum height=1.3cm
  },
  flecha/.style={
    -{Latex[length=1.8mm,width=1.4mm]},
    draw=black!65, line width=0.5pt
  },
  rotulo/.style={
    rotate=90, font=\tiny\bfseries, text=black!70, align=center
  }
]

%% ------------------------- Efectos indirectos ------------------------
\node[nodo] (I1) at (-8.4,0) {Pérdida potencial de clientes y oportunidades de venta};
\node[nodo] (I2) at (-2.8,0) {Disminución de la confianza en el servicio ofrecido};
\node[nodo] (I3) at ( 2.8,0) {Mayores costos de coordinación y corrección};
\node[nodo] (I4) at ( 8.4,0) {Dificultades en la planificación de salidas turísticas};

%% -------------------------- Efectos directos -------------------------
\node[nodo] (D1) at (-8.4,-2.6) {Demoras en la confirmación de reservas};
\node[nodo] (D2) at (-2.8,-2.6) {Errores en las reservas, cupos duplicados o mal asignados};
\node[nodo] (D3) at ( 2.8,-2.6) {Información contradictoria sobre la disponibilidad};
\node[nodo] (D4) at ( 8.4,-2.6) {Reprogramaciones, reclamos y cancelaciones};

%% --------------------------- Problema central ------------------------
\node[central] (C) at (0,-5.2) {Ineficiencia en el proceso de venta de paquetes turísticos en Vizcacha Tours, Potosí};

%% --------------------------- Causas directas -------------------------
\node[nodo] (CD1) at (-8.4,-7.8) {Verificación no uniforme de cupos antes de confirmar una reserva};
\node[nodo] (CD2) at (-2.8,-7.8) {Desfase en la actualización de ventas entre distintos canales de atención};
\node[nodo] (CD3) at ( 2.8,-7.8) {Cambios de fechas o capacidad que no se reflejan oportunamente en las reservas};
\node[nodo] (CD4) at ( 8.4,-7.8) {Criterios poco claros para modificar, cancelar y liberar reservas};

%% ------------------------- Causas indirectas -------------------------
\node[nodo] (CI1) at (-8.4,-10.4) {Responsabilidades de verificación y actualización sin definición suficiente};
\node[nodo] (CI2) at (-2.8,-10.4) {Información de salidas y reservas sin una fuente compartida y actualizada};
\node[nodo] (CI3) at ( 2.8,-10.4) {Falta de controles para identificar cambios y prevenir conflictos de capacidad};
\node[nodo] (CI4) at ( 8.4,-10.4) {Procedimientos de seguimiento y validación insuficientemente estandarizados};

%% ------------------------------ Flechas ------------------------------
%% Efecto directo -> efecto indirecto
\foreach \d/\i in {D1/I1,D2/I2,D3/I3,D4/I4} {
  \draw[flecha] (\d.north) -- (\i.south);
}

%% Problema central -> efectos directos
\foreach \d/\x in {D1/-8.4,D2/-2.8,D3/2.8,D4/8.4} {
  \draw[flecha] (C.north -| \x,0) -- (\d.south);
}

%% Causas directas -> problema central
\foreach \c/\x in {CD1/-8.4,CD2/-2.8,CD3/2.8,CD4/8.4} {
  \draw[flecha] (\c.north) -- (C.south -| \x,0);
}

%% Causas indirectas -> causas directas
\foreach \c/\d in {CI1/CD1,CI2/CD2,CI3/CD3,CI4/CD4} {
  \draw[flecha] (\c.north) -- (\d.south);
}

%% --------------------------- Rótulos de nivel -------------------------
\node[rotulo] at (-10.95,0)     {E. INDIRECTOS};
\node[rotulo] at (-10.95,-2.6)  {E. DIRECTOS};
\node[rotulo] at (-10.95,-5.2)  {PROBLEMA};
\node[rotulo] at (-10.95,-7.8)  {C. DIRECTAS};
\node[rotulo] at (-10.95,-10.4) {C. INDIRECTAS};

\end{tikzpicture}
}
\vspace{2pt}
\begin{minipage}{0.88\linewidth}
  \nota{Fuente: elaboración propia a partir del diagnóstico preliminar.}
\end{minipage}
\end{figuraapa}
\end{landscape}
